\documentclass[11pt]{article}
\usepackage[margin=1in]{geometry}
\usepackage[T1]{fontenc}
\usepackage{lmodern}
\usepackage{amsmath,amssymb}
\usepackage[hidelinks]{hyperref}

\title{Literature Status for arXiv:2412.07920:\\
Sharp Spectral Multipliers Beyond M\'etivier Groups}
\author{Run fa\_banach\_001, agent\_lane\_05}
\date{11 August 2026}

\begin{document}
\maketitle

\section*{Status}

This is a \texttt{literature\_already\_answered} packet.  In
\emph{Spectral multipliers on M\'etivier groups}, Niedorf asks whether
$p$-specific spectral multiplier estimates with the sharp threshold
\[
  s>d\left|\frac1p-\frac12\right|
\]
hold beyond the class of M\'etivier groups \cite{NiedorfM}.  His later paper
\emph{Spectral multipliers on two-step stratified Lie groups with degenerate
group structure} gives an affirmative answer for explicit, genuinely
non-M\'etivier subclasses \cite{NiedorfD}.

\section*{Source Question}

Let $G$ be a two-step stratified Lie group of topological dimension $d$, let
$L$ be a sub-Laplacian, and let
\[
 J_\mu\colon\mathfrak g_1\longrightarrow\mathfrak g_1,
 \qquad
 \langle J_\mu x,x'\rangle=\mu([x,x']).
\]
A M\'etivier group is characterized by the invertibility of $J_\mu$ for every
$\mu\ne0$.  After proving the sharp estimate in that nondegenerate setting,
arXiv:2412.07920 asks whether the same $p$-specific regularity threshold holds
beyond M\'etivier groups.  It points to weighted Plancherel estimates as the
main obstruction when $J_\mu$ has a nontrivial kernel.

\section*{Later Answer}

The introduction of arXiv:2501.16262 explicitly cites the M\'etivier paper,
describes the same question for groups where nondegeneracy fails, and then
turns to a class in which $J_\mu$ may have a nontrivial kernel.  Its
Theorem~1.1 proves the following.  Suppose the first layer has an orthogonal
block decomposition satisfying the paper's Assumptions A and B, and write
$\bar d_1$ for the total rank of the nonzero spectral blocks.  If
$\bar d_1\ge d_2-1$, then, in its stated range
$1\le p\le p_{\bar d_1,d_2}$, every bounded Borel function $F$ satisfying
\[
 \|F\|_{L^2_{s,\mathrm{sloc}}}<\infty,
 \qquad
 s>d\left(\frac1p-\frac12\right),
\]
gives an $L^p(G)$-bounded operator $F(L)$ for $p>1$.  The corresponding
Bochner--Riesz estimate also holds.  The supporting paper states explicitly
that its theorem is completely new in the non-M\'etivier case.

The assumptions are verified there for concrete degenerate families.  One is
a central product of $N$ copies of the free two-step group $N_{3,2}$; for
$N\ge2$ the multiplier conclusion has a nonempty interval with $p>1$.
Another is the Heisenberg--Reiter family $\mathbb H_{N,d_2}$ under the stated
dimensional condition.  In both families the displayed kernels of $J_\mu$
show directly that the relevant examples are not M\'etivier.  Hence sharp
$p$-specific multiplier estimates do hold beyond the M\'etivier class.

\section*{Scope}

The later theorem does not cover every two-step stratified Lie group.  It
requires two structural assumptions and a dimensional inequality; its own
discussion says that the needed Plancherel estimates are unavailable in
general.  Thus the existential ``beyond M\'etivier'' question has a rigorous
affirmative answer for substantial degenerate subclasses, while the universal
all-two-step problem remains open.

\section*{Provenance Check}

The identification is explicit rather than merely inferential:
arXiv:2501.16262 cites arXiv:2412.07920 as its M\'etivier predecessor, restates
the question caused by failure of nondegeneracy, and labels its main theorem
new in the non-M\'etivier setting.  Both version-2 PDFs are included in this
packet.

\begin{thebibliography}{9}

\bibitem{NiedorfM}
L.~Niedorf,
\emph{Spectral multipliers on M\'etivier groups},
arXiv:2412.07920v2 (2025),
\href{https://arxiv.org/abs/2412.07920}{arXiv:2412.07920}.

\bibitem{NiedorfD}
L.~Niedorf,
\emph{Spectral multipliers on two-step stratified Lie groups with degenerate
group structure},
arXiv:2501.16262v2 (2025),
\href{https://arxiv.org/abs/2501.16262}{arXiv:2501.16262}.

\end{thebibliography}

\end{document}
